% SPDX-License-Identifier: CC-BY-SA-4.0
% Part: Introduction
% Section: Words and languages
% Exercise: Words

\begingroup

\begin{exercice}[Mots]\label{exo:introduction/languages/words}

  Soit $x = abbcc$ un mot sur l'alphabet $\Sigma = \{a, b, c\}$.
  \begin{question}
  \item Quelle est la valeur de $|x|$ ?
  \item Donner un mot de $\Sigma^3$ qui n'est pas un facteur de $x$.
  \item Donner un sous-mot\footnote{c'est-à-dire une partie de la liste des symboles de $x$ pas forcément contiguë} de $x$ qui n'est pas un facteur de $x$.
  \item Donner tous les facteurs de $x$ qui appartiennent à $\Sigma^3$ .
  \item Donner l'ensemble $\mathit{Pref}(x)$ des préfixes de $x$.
  \item Donner l'ensemble $\mathit{SuffProp}(x)$ des suffixes propres de $x$.
  \end{question}
  Reprendre ces questions en considérant maintenant $x$ comme un mot sur
  l'alphabet $\Sigma = \{ab, ac, bc, c\}$.

\end{exercice}

\begin{correction}

  \begin{itemize}
  \item Sur $\Sigma = \{a, b, c\}$, $x = abbcc = a\cdot b\cdot b\cdot c\cdot c$ :
    \begin{question}
    \item $|x| = 5$
    \item par exemple $bac$
    \item par exemple $acc$
    \item $\{abb, bbc, bcc\}$
    \item $\mathit{Pref}(x) = \{\varepsilon, a, ab, abb, abbc, abbcc\}$
    \item $\mathit{SuffProp}(x) = \{\varepsilon, c, cc, bcc, bbcc\}$ (tous les suffixes sauf le mot entier)
    \end{question}
  \item Sur $\Sigma = \{ab, ac, bc, c\}$, $x = abbcc = ab\cdot bc\cdot c$ :
    \resetquestion
    \begin{question}
    \item $|x| = 3$
    \item par exemple $ab\cdot ab\cdot ab$
    \item il n'y en a qu'un : $ab\cdot c$
    \item $\{ab.bc.c\}$
    \item $\mathit{Pref}(x) = \{\varepsilon, ab, ab\cdot bc, ab\cdot bc\cdot c\}$
    \item $\mathit{SuffProp}(x) = \{\varepsilon, c, bc\cdot c\}$ (tous les suffixes sauf le mot entier)
    \end{question}
  \end{itemize}

\end{correction}

\endgroup
\endinput
